\documentclass[12pt, a4paper]{article}
\usepackage[UTF8, zihao=-4]{ctex}
\usepackage[top=2.5cm, bottom=2.5cm, left=2.8cm, right=2.8cm]{geometry}
\usepackage{fontspec}
\setmainfont{Times New Roman}
\usepackage{setspace}
\setstretch{1.5}
\usepackage{booktabs}
\usepackage{array}
\usepackage{makecell}
\usepackage{multirow}
\usepackage{adjustbox}
\usepackage{amsmath, amssymb}
\usepackage{float}
\usepackage{caption}
\usepackage{fancyhdr}
\setlength{\headheight}{15pt}
\pagestyle{fancy}
\fancyhf{}
\fancyhead[C]{\small 实验数据记录单}
\fancyfoot[C]{\small\thepage}

\begin{document}

\begin{center}
  {\heiti\zihao{3} 刀口阴影法测量光学系统像差综合实验 \\ 实验数据记录单} \\[1cm]
  \begin{tabular}{llll}
    姓\quad 名：& \underline{\hspace{4cm}} &
    学\quad 号：& \underline{\hspace{4cm}} \\[0.5cm]
    班\quad 级：& \underline{\hspace{4cm}} &
    实验日期：  & \underline{\hspace{4cm}} \\[0.5cm]
    指导教师：  & \underline{\hspace{4cm}} &
    教师签名：  & \underline{\hspace{4cm}} \\
  \end{tabular}
\end{center}
\vspace{1cm}

\subsection*{表 1 \quad 红色光源轴向球差（环带光测试法）}
\begin{table}[H]
  \centering
  \begin{adjustbox}{max width=\textwidth}
    \begin{tabular}{cccc}
      \toprule
      次数 & $X_1$/mm & $X_2$/mm & 轴向球差 $\Delta X=X_2-X_1$/mm \\
      \midrule
      1 & & & \\
      2 & & & \\
      3 & & & \\
      \midrule
      平均值 & & & \\
      \bottomrule
    \end{tabular}
  \end{adjustbox}
\end{table}
{\footnotesize 记录说明：$X_1$ 为最小环带光阑对应的会聚位置读数，$X_2$ 为中号环带光阑对应的会聚位置读数。读数单位统一为 mm。}

\vspace{0.8cm}

\subsection*{表 2 \quad 透镜像散数据记录}
\begin{table}[H]
  \centering
  \begin{adjustbox}{max width=\textwidth}
    \begin{tabular}{cccc}
      \toprule
      旋转角度 & $X_1$/mm & $X_2$/mm & 透镜像散 $\Delta X=X_2-X_1$/mm \\
      \midrule
      10$^\circ$ & & & \\
      20$^\circ$ & & & \\
      30$^\circ$ & & & \\
      \midrule
      平均值 & & & \\
      \bottomrule
    \end{tabular}
  \end{adjustbox}
\end{table}
{\footnotesize 记录说明：$X_1$、$X_2$ 分别对应两条互相垂直焦线的轴向位置。调节螺旋丝杆时不要超出其调节范围。}

\vspace{1cm}
\noindent\textbf{简要现象记录：}

\vspace{2cm}

\noindent\textbf{误差来源记录：}

\vspace{2cm}

\end{document}
